% Source - https://tex.stackexchange.com/a/457624
% Posted by user121799
% Retrieved 2026-03-17, License - CC BY-SA 4.0

\documentclass[tikz,border=8mm]{standalone}
\usetikzlibrary{shapes.geometric,shapes.symbols,fit,positioning,shadows}
% https://tex.stackexchange.com/a/12039/121799
\makeatletter
% https://tex.stackexchange.com/a/103691/121799
\pgfdeclareshape{document}{
\inheritsavedanchors[from=rectangle] % this is nearly a rectangle
\inheritanchorborder[from=rectangle]
\inheritanchor[from=rectangle]{center}
\inheritanchor[from=rectangle]{north}
\inheritanchor[from=rectangle]{north east}
\inheritanchor[from=rectangle]{north west}
\inheritanchor[from=rectangle]{south}
\inheritanchor[from=rectangle]{south east}
\inheritanchor[from=rectangle]{south west}
\inheritanchor[from=rectangle]{west}
\inheritanchor[from=rectangle]{east}
\backgroundpath{%
\southwest \pgf@xa=\pgf@x \pgf@ya=\pgf@y
\northeast \pgf@xb=\pgf@x \pgf@yb=\pgf@y
\pgf@xc=\pgf@xb \advance\pgf@xc by-5pt % this should be a parameter
\pgf@yc=\pgf@ya \advance\pgf@yc by5pt
\pgfpathmoveto{\pgfpoint{\pgf@xa}{\pgf@ya}}
\pgfpathlineto{\pgfpoint{\pgf@xa}{\pgf@yb}}
\pgfpathlineto{\pgfpoint{\pgf@xb}{\pgf@yb}}
\pgfpathlineto{\pgfpoint{\pgf@xb}{\pgf@yc}}
\pgfpathlineto{\pgfpoint{\pgf@xc}{\pgf@ya}}
\pgfpathclose
% add little corner
\pgfpathmoveto{\pgfpoint{\pgf@xc}{\pgf@ya}}
\pgfpathlineto{\pgfpoint{\pgf@xc}{\pgf@yc}}
\pgfpathlineto{\pgfpoint{\pgf@xb}{\pgf@yc}}
\pgfpathclose
}
}
\makeatother

\begin{document}

\tikzset{doc/.style={document,fill=blue!10,draw,thin,minimum
height=1.2cm,align=center},
pics/.cd,
pack/.style={code={%
\draw[fill=blue!50,opacity=0.2] (0,0) -- (0.5,-0.25) -- (0.5,0.25) -- (0,0.5) -- cycle;
\draw[fill=blue!50,opacity=0.2] (0,0) -- (-0.5,-0.25) -- (-0.5,0.25) -- (0,0.5) -- cycle;
\draw[fill=blue!60,opacity=0.2] (0,0) -- (-0.5,-0.25) -- (0,-0.5) -- (0.5,-0.25) -- cycle;
\draw[fill=blue!60] (0,0) -- (0.25,0.125) -- (0,0.25) -- (-0.25,0.125) -- cycle;
\draw[fill=blue!50] (0,0) -- (0.25,0.125) -- (0.25,-0.125) -- (0,-0.25) -- cycle;
\draw[fill=blue!50] (0,0) -- (-0.25,0.125) -- (-0.25,-0.125) -- (0,-0.25) -- cycle;
\draw[fill=blue!50,opacity=0.2] (0,-0.5) -- (0.5,-0.25) -- (0.5,0.25) -- (0,0) -- cycle;
 \draw[fill=blue!50,opacity=0.2] (0,-0.5) -- (-0.5,-0.25) -- (-0.5,0.25) -- (0,0) -- cycle;
\draw[fill=blue!60,opacity=0.2] (0,0.5) -- (-0.5,0.25) -- (0,0) -- (0.5,0.25) -- cycle;
}}}

\begin{tikzpicture}[
    font=\sffamily,
    every label/.append style={font=\small\sffamily,align=center}]

% hardware layer nodes
\node[doc,fill=red!10]
    (Distance)
    {Ultrasonic\\Distance\\(5V)};
\node[right=2cm of Distance,doc,fill=blue!10]
    (Receiver)
    {Break Beam\\Receiver};
\node[right=2cm of Receiver,doc,fill=blue!10]
    (Emitter)
    {Break Beam\\Emitter};
\node[right=2cm of Emitter,doc,fill=blue!10]
    (LEDs)
    {LEDs\\};
\node[right=2cm of LEDs,doc,fill=blue!10]
    (Buzzer)
    {Passive\\Buzzer};

% BBB pin layer nodes
\node[below=4cm of Distance,doc,fill=blue!10]
    (Ground)
    {Ground};
\node[below=4cm of Receiver,doc,fill=blue!10]
    (Power)
    {Power};
\node[below=4cm of Emitter,doc,fill=blue!10]
    (GPIO_In)
    {GPIO In};
\node[below=4cm of LEDs,doc,fill=blue!10]
    (GPIO_Out)
    {GPIO Out};
\node[below=4cm of Buzzer,doc,fill=blue!10]
    (PWM)
    {PWM};

% BBB layer
\node[below=2cm of GPIO_In,doc,fill=blue!10]
    (BBB)
    {BBB};

% Host layer
\node[below=2cm of BBB,doc,fill=blue!10]
    (Host)
    {Host};

% links between layers 1 and 2
\draw (Distance) -- (Ground);
\draw (Distance) -- (Power);
\draw (Distance) -- (GPIO_In);
\draw (Distance) -- (GPIO_Out);
\draw (Receiver) -- (Ground);
\draw (Receiver) -- (Power);
\draw (Receiver) -- (GPIO_In);
\draw (Emitter) -- (Ground);
\draw (Emitter) -- (Power);
\draw (LEDs) -- (Ground);
\draw (LEDs) -- (GPIO_Out);
\draw (Buzzer) -- (Ground);
\draw (Buzzer) -- (PWM);
\draw (BBB) -- (Ground);
\draw (BBB) -- (Power);
\draw (BBB) -- (GPIO_In);
\draw (BBB) -- (GPIO_Out);
\draw (BBB) -- (PWM);
\draw (BBB) -- node[right]{USB}(Host);


\node[draw,dashed,rounded corners,
    fit=(Distance)(Receiver)(Emitter),
    inner sep=10pt,label={left:{Sensors}}]
    (fit1)
    {};
\node[draw,dashed,rounded corners,
    fit=(LEDs)(Buzzer),
    inner sep=10pt,label={left:{Outputs}}]
    (fit2)
    {};
\node[draw,dashed,rounded corners,
    fit=(Ground)(Power)(GPIO_In)(GPIO_Out)(PWM),
    inner sep=10pt,label={left:{BBB\\Pins}}]
    (fit3)
    {};

\end{tikzpicture}
\end{document}
